\section{Mapa de Prioridad}

Este documento resume los ejercicios de las guias de estudio que tienen mayor retorno para el examen: no son todos los ejercicios, sino los que entrenan algoritmos que pueden aparecer con numeros cambiados.

\begin{center}
\begin{tabularx}{\linewidth}{p{0.16\linewidth}X X p{0.12\linewidth}}
\toprule
Guia & Ejercicio & Por que importa & Prioridad \\
\midrule
Guia 11 & Muestreo de aceptacion binomial & Traduce AQL/LTPD y riesgos a probabilidades de aceptar o rechazar lotes. & Alta \\
Guia 11 & Grafico $\bar X$ con $\sigma$ conocida & Es el caso mas rapido de limites de control. & Alta \\
Guia 10 & Kanban + control $\bar X/R$ & Cruza JIT, calidad y diagnostico de proceso. & Alta \\
Guia 9 & Profundidad optima de pallets & Ejercicio logistico corto y muy formulable. & Alta \\
Guia 8 & M/M/1, variabilidad y fallas & Refuerza Little, utilizacion, CV y disponibilidad. & Media \\
\bottomrule
\end{tabularx}
\end{center}

\begin{alerta}{como estudiar esto}
El examen probablemente no copia el enunciado, pero si copia la estructura: te cambian demanda, nivel de confianza, defectos, ciclos o costos. Por eso cada bloque define variables y deja el algoritmo limpio.
\end{alerta}

La lectura recomendada es:
\begin{enumerate}
  \item Calidad: dominar primero muestreo de aceptacion y limites $\bar X/R$.
  \item Kanban/JIT: entender que mas defectos obligan a producir mas unidades brutas para cumplir demanda buena.
  \item Bodegas: memorizar la formula de profundidad y sus supuestos.
  \item Variabilidad/colas: saber identificar $\lambda$, $\mu$, $\rho$, $L$, $W$ y disponibilidad.
\end{enumerate}
